\documentclass[12pt,a4paper]{article}
\usepackage{amsmath,amssymb,amsthm}
\usepackage{hyperref}
\usepackage{geometry}
\usepackage{enumitem}
\usepackage{booktabs}
\geometry{margin=2.5cm}

\newtheorem{theorem}{Theorem}[section]
\newtheorem{lemma}[theorem]{Lemma}
\newtheorem{corollary}[theorem]{Corollary}
\newtheorem{proposition}[theorem]{Proposition}
\theoremstyle{definition}
\newtheorem{definition}[theorem]{Definition}
\theoremstyle{remark}
\newtheorem{remark}[theorem]{Remark}

\title{Grand Unification in Recognition Science:\\
One Ledger, Three Sectors}

\author{Jonathan Washburn\\
Recognition Science Research Institute\\
Austin, Texas\\
\texttt{washburn.jonathan@gmail.com}}

\date{March 2026}

\begin{document}
\maketitle

\begin{abstract}
We present the Recognition Science approach to grand unification.
Unlike conventional GUTs (which embed
$SU(3) \times SU(2) \times U(1)$ into a larger group like $SU(5)$
or $SO(10)$), RS unification is structural: all three gauge sectors
originate from a single recognition ledger.  The three sectors are
different topological projections of the same $D$-cube ($D = 3$):
color from face-pairs, weak isospin from spinors, hypercharge from
overall phase.  This explains why the couplings are related without
requiring a GUT scale or predicting proton decay.  The ``unification
scale'' is not a high-energy threshold but the ledger itself---already
unified at all scales.  All structural results follow from the RS forcing chain~\cite{RS_Axioms}.
\end{abstract}

%% ============================================================
\section{Introduction}
\label{sec:intro}
%% ============================================================

Grand unification attempts to explain why the three gauge couplings
appear to converge at high energies~\cite{GeorgiGlashow1974}.
Georgi and Glashow~\cite{GeorgiGlashow1974} proposed $SU(5)$ as the
minimal GUT: fermions in $\mathbf{5}^* \oplus \mathbf{10}$, with
$M_{\mathrm{GUT}} \sim 10^{16}$~GeV from RG running.  Pati and
Salam~\cite{PatiSalam1974} extended via $SU(4) \times SU(2)_L \times
SU(2)_R$; Georgi~\cite{Georgi1975} and Fritzsch--Minkowski~\cite{FritzschMinkowski1975}
proposed $SO(10)$ with a $\mathbf{16}$ spinor.  All require a GUT
scale and predict proton decay and monopoles.

\textbf{RS alternative.}
Recognition Science~\cite{RS_Axioms} (RS) offers a different kind of unification.
The cost functional $J(x) = \tfrac{1}{2}(x + x^{-1}) - 1$ and the
golden ratio $\varphi = (1+\sqrt{5})/2$ force $D = 3$ spatial
dimensions.  The $D$-cube geometry then forces
$SU(3) \times SU(2) \times U(1)$: one ledger, three sectors as
different projections.  No GUT scale is needed---the structure is
already unified at all scales.  No proton decay or monopoles are
predicted.

%% ============================================================
\section{Recognition Science Framework}
\label{sec:framework}
%% ============================================================

Recognition Science derives all physics from a single primitive,
the \emph{recognition cost functional} $J(x)$, uniquely determined
by three axioms.

\begin{definition}[RS Cost Functional]
For $x > 0$:
\begin{equation}
  J(x) = \frac{1}{2}\!\left(x + x^{-1}\right) - 1.
\end{equation}
\end{definition}

\noindent\textbf{Axiom A1 (Normalization):} $J(1) = 0$.\\
\textbf{Axiom A2 (Recognition Composition Law, RCL):}
$J(xy) + J(x/y) = 2J(x)J(y) + 2J(x) + 2J(y)$ for all $x, y > 0$.\\
\textbf{Axiom A3 (Calibration):} $J''_{\log}(0) = 1$, where
$J_{\log}(t) := J(e^t)$.

\begin{theorem}[Cost Uniqueness]
\label{thm:unique}
The unique continuous solution to \textup{(A1)--(A3)} is
$J(x) = \tfrac{1}{2}(x + x^{-1}) - 1$.
\end{theorem}

\begin{proof}[Proof sketch]
Setting $f(t) = J_{\log}(t) + 1$ and rewriting \textup{(A2)} gives
the d'Alembert functional equation $f(s+t) + f(s-t) = 2f(s)f(t)$.
By the Acz\'el classification theorem~\cite{Aczel1966}, continuous
solutions with $f(0) = 1$ and $f''(0) = 1$ are $f(t) = \cosh t$,
yielding $J(x) = \tfrac{1}{2}(x + x^{-1}) - 1$.
\end{proof}

\noindent The golden ratio $\varphi = (1+\sqrt{5})/2$ is forced by
self-similarity (the $\varphi$-ladder).  Spatial dimension $D = 3$
is forced by the gap-45 synchronization $\operatorname{lcm}(8, 45)
= 360$.  The $D = 3$ cube structure is the geometric foundation of
all three gauge sectors.

%% ============================================================
\section{Structural Unification}
\label{sec:structural}
%% ============================================================

\begin{proposition}[Single Ledger --- Structural Argument]
\label{thm:single-ledger}
All three gauge sectors (strong, weak, electromagnetic) originate
from a single recognition ledger.  The apparent multiplicity is
a projection effect.
\end{proposition}

\begin{proof}[Structural argument]
In RS, the entire physical state is encoded in the recognition ledger
$L = \{(x_i, J(x_i))\}$, a single double-entry accounting system.
The three gauge symmetries are not separate structures imposed on the
ledger but are consequences of the three independent structural
features of the $D=3$ cube (Theorem~\ref{thm:projections}):
face-pairs, spinor lifts, and $J$-cost phase.  Since all three derive
from the same cube geometry and cost functional, they share one
common origin.  There is no separate ``strong sector ledger'' or
``weak sector ledger''---only the three projections of one structure.
\end{proof}

\begin{proposition}[Common Origin from $D = 3$ --- Structural Argument]
\label{thm:common-origin}
$N_c = 3$ (strong), $N_{\mathrm{gen}} = 3$ (weak generations), and the $U(1)$
hypercharge all follow from the single structural input $D = 3$.
\end{proposition}

\begin{proof}[Structural argument]
From Lemma~\ref{lem:face-pairs}: $\binom{3}{2} = 3$ independent
coordinate planes $\Rightarrow N_c = 3$ color channels ($SU(3)$).  The $SO(3)$ spinor
representation is a doublet of $SU(2)$, giving two weak isospin states.
Three colors times two isospin states times the overall $J$-cost phase
yields one generation.  The $\varphi$-ladder sectors before first
periodicity of the recognition angle $\theta_0 = \arccos(1/4)$ give
three distinct sectors, motivating $N_{\mathrm{gen}} = 3$; however, the
rigorous derivation of exactly three generations is an identified open
problem in RS (see Remark below).
\end{proof}

\begin{remark}[Open problem: three generations]
\label{rem:3gen}
The argument that $\theta_0 = \arccos(1/4)$ forces exactly three
$\varphi$-rung sectors before periodicity is a structural motivation,
not a rigorous derivation.  A complete proof would require showing that
the $\varphi$-ladder restricted to the recognition angle has exactly
three periodic sectors and that each sector corresponds to a physical
generation of fermions.  This is one of the most important open problems
in the RS framework.
\end{remark}

%% ============================================================
\section{Three Sectors as Projections of One $D$-Cube}
\label{sec:projections}
%% ============================================================

\begin{definition}[$D$-Cube]
The $D$-cube has vertices $\{0,1\}^D$.  For $D = 3$: 8 vertices, 12
edges, 6 faces (opposite-face pairs: 3), and 3 independent coordinate
planes (2-plane pairs: $\binom{3}{2} = 3$).  For $D = 3$ specifically,
the number of opposite-face pairs equals the number of 2-plane pairs
($D = \binom{D}{2} = 3$); this coincidence holds only at $D = 3$.
\end{definition}

\begin{lemma}[2-Plane Count for $D = 3$]
\label{lem:face-pairs}
The $D = 3$ cube has $\binom{3}{2} = 3$ coordinate planes:
$(xy)$, $(xz)$, $(yz)$.  These define 3 independent complex
recognition channels.
\end{lemma}

\begin{proposition}[Projection Structure --- Structural Argument]
\label{thm:projections}
The three gauge sectors are structural projections of the $D = 3$ cube:
(i) \textbf{Color ($SU(3)$)}: the 3 coordinate planes each index a
complex recognition channel, motivating $SU(3)$;
(ii) \textbf{Weak ($SU(2)$)}: the $SO(3)$ rotation group of the 3-cube
has a doublet spinor representation lifting to $SU(2)$;
(iii) \textbf{Hypercharge ($U(1)$)}: $J(x) = \tfrac{1}{2}(x + x^{-1}) - 1$
is invariant under phase rotation $x \mapsto e^{i\theta}x$, giving $U(1)$.
\end{proposition}

\begin{remark}[Status of the gauge group derivation]
\label{rem:gauge-status}
The identifications (i)--(iii) are structural: the face-pair count 3
corresponds to $SU(3)$; the $SO(3)$ spinor lift gives $SU(2)$; the
$J$-phase symmetry gives $U(1)$.  The gap in each case is showing that
the \emph{representation content} (which irreducible representations the
quarks and leptons transform under) follows from the RS ledger without
additional assumptions.  This is an identified open problem.
\end{remark}

\begin{corollary}
\label{cor:no-embedding}
RS unification does not embed $SU(3) \times SU(2) \times U(1)$
into a larger group.  The three sectors are already facets of
one structure.
\end{corollary}

%% ============================================================
\section{Gauge Coupling Relations Without Running}
\label{sec:couplings}
%% ============================================================

Conventional GUTs relate couplings by renormalization-group running
to a common scale $M_{\mathrm{GUT}}$.  RS relates them by geometry
at all scales.

\begin{proposition}[Coupling Relations from $D = 3$ --- Summary]
\label{thm:couplings}
The gauge couplings are determined by the $D$-cube structure with zero
free parameters:
\begin{itemize}[noitemsep]
\item $\sin^2\theta_W = (3 - \varphi)/6 \approx 0.230$ from the
  electroweak $\varphi$-structure;
\item $\alpha_s = 2/17 \approx 0.1176$ at $M_Z$ from the 17
  2D crystallographic (wallpaper) groups;
\item $\alpha^{-1} \approx 137$ from the cube passive-edge count
  and $f_{\mathrm{gap}}$ correction.
\end{itemize}
Each derivation is developed in full in companion RS papers on
the Weinberg angle, strong coupling, and fine-structure constant,
respectively~\cite{RS_Axioms}.
\end{proposition}

\begin{remark}[Scale independence and falsifiability]
The coupling values are set geometrically by $D = 3$ and hold at all
scales without RG running.  This is a sharp departure from conventional
GUTs: the RS values (e.g.\ $\sin^2\theta_W = (3-\varphi)/6 \approx
0.2306$) should be \emph{scale-independent}, while measured couplings
run with energy.  If precision measurements establish significant
running in $\sin^2\theta_W$ below $M_Z$, the RS geometric derivation
would be under pressure.
\end{remark}

%% ============================================================
\section{No Proton Decay}
\label{sec:proton-decay}
%% ============================================================

The Georgi--Glashow $SU(5)$ GUT predicts proton decay via $X$/$Y$
leptoquark bosons mediating $q \leftrightarrow \ell$, e.g.\
$p \to e^+ \pi^0$.  Predicted $\tau_p \sim 10^{30}$--$10^{35}$~yr;
experiments give $\tau_p > 10^{34}$~yr, ruling out minimal $SU(5)$.

\begin{theorem}[No Proton Decay in RS]
\label{thm:no-proton-decay}
RS does not predict proton decay.
\end{theorem}

\begin{proof}
No embedding into $SU(5)$ or $SO(10)$ $\Rightarrow$ no $X$/$Y$
bosons.  The three sectors are facets of one ledger, not broken
remnants.  No gauge bosons mediate $q \leftrightarrow \ell$.
\end{proof}

%% ============================================================
\section{No GUT Monopoles}
\label{sec:monopoles}
%% ============================================================

When $SU(5)$ breaks to $SU(3) \times SU(2) \times U(1)$, the vacuum
manifold has $\pi_2(G/H) \neq 0$, so monopoles are topologically
required (t'Hooft--Polyakov~\cite{tHooft1974,Polyakov1974}), with
mass $\sim M_{\mathrm{GUT}}/\alpha \sim 10^{16}$~GeV.

\begin{theorem}[No GUT Monopoles]
\label{thm:no-monopoles}
RS does not predict grand-unification monopoles.
\end{theorem}

\begin{proof}
Monopoles require $G \to H$ with $\pi_2(G/H) \neq 0$.  In RS there
is no such breaking; $SU(3) \times SU(2) \times U(1)$ is the direct
product of projections, not a broken remnant.  $\pi_2$ is trivial.
\end{proof}

%% ============================================================
\section{Comparison with Conventional GUTs}
\label{sec:comparison}
%% ============================================================

\begin{center}
\renewcommand{\arraystretch}{1.2}
\begin{tabular}{lcc}
\toprule
& \textbf{Conventional GUT} & \textbf{RS} \\
\midrule
Unification & Larger group ($SU(5)$, $SO(10)$) & Single ledger \\
GUT scale & $\sim 10^{16}$~GeV (RG running) & None (all scales) \\
Proton decay & Predicted ($X$/$Y$ bosons) & Not predicted \\
Monopoles & Topologically required & Not predicted \\
Couplings & Run to $M_{\mathrm{GUT}}$ & From $D$-cube geometry \\
Parameters & $M_{\mathrm{GUT}}$, breaking sector & Zero \\
\bottomrule
\end{tabular}
\end{center}

%% ============================================================
\section{Predictions and Falsifiers}
\label{sec:predictions}
%% ============================================================

\begin{enumerate}[label=\textbf{P\arabic*}.]
\item \textbf{P1}: No proton decay.  Falsifier: $p \to e^+ \pi^0$
  or similar.
\item \textbf{P2}: No GUT monopoles.  Falsifier: monopoles
  $\sim 10^{16}$~GeV.
\item \textbf{P3}: $\sin^2\theta_W = (3-\varphi)/6$, $\alpha_s =
  2/17$, $\alpha^{-1} \approx 137$.  Falsifier: measurements outside
  bands.
\item \textbf{P4}: Exactly three generations.  Falsifier: fourth
  light generation.
\item \textbf{P5}: No $X$/$Y$ bosons.  Falsifier: $q \leftrightarrow
  \ell$ via new gauge bosons.
\end{enumerate}

%% ============================================================
\section{Conclusion}
\label{sec:conclusion}
%% ============================================================

Grand unification in RS is structural: one ledger yields three
gauge sectors as different projections of the $D$-cube.  No GUT
scale, no proton decay, no monopoles.  The coupling relations
derive from $D = 3$ geometry, not from renormalization-group
running.  RS offers a falsifiable alternative to conventional
GUTs, with distinct experimental signatures.

\begin{thebibliography}{99}

\bibitem{Aczel1966}
J.~Acz\'el, \emph{Lectures on Functional Equations and Their Applications},
Academic Press, New York (1966).

\bibitem{GeorgiGlashow1974}
H.~Georgi and S.~L.~Glashow, Phys.\ Rev.\ Lett.\ \textbf{32}, 438 (1974).

\bibitem{Georgi1975}
H.~Georgi, in \emph{Particles and Fields} (AIP, New York, 1975), p.~575.

\bibitem{PatiSalam1974}
J.~C.~Pati and A.~Salam, Phys.\ Rev.\ D \textbf{10}, 275 (1974).

\bibitem{FritzschMinkowski1975}
H.~Fritzsch and P.~Minkowski, Ann.\ Phys.\ \textbf{93}, 193 (1975).

\bibitem{tHooft1974}
G.~'t Hooft, Nucl.\ Phys.\ B \textbf{79}, 276 (1974).

\bibitem{Polyakov1974}
A.~M.~Polyakov, JETP Lett.\ \textbf{20}, 194 (1974).

\bibitem{RS_Axioms}
J.~Washburn, ``The Algebra of Reality: A Recognition Science Derivation
of Physical Law,'' \emph{Axioms} \textbf{15}(2), 90 (2026).
\texttt{doi:10.3390/axioms15020090}
\end{thebibliography}

\end{document}
